\section{Strojové učení a~algoritmy zpětnovazebního učení}

Strojové učení je samostatná oblast umělé inteligence, která se zabývá vývojem metod umožňujících
počítačovým systémům automaticky zlepšovat své chování na základě zkušeností bez nutnosti explicitního
programování všech rozhodovacích pravidel. Takové systémy se učí přímo z~příkladů nebo interakcí
s~prostředím, identifikují vzory a~vytvářejí modely, které generalizují i~na nové situace, jež
nebyly přímo obsaženy v~tréninkových datech.

Takový přístup je zásadní zejména tam, kde ruční definice pravidel nevede k~dobrému výsledku – například
při řízení autonomních agentů v~komplexních, nelineárních a~stochastických prostředích nebo při
adaptivním rozhodování v~dynamických situacích, které nelze přesně popsat jednou sadou pravidel.

Obecně rozlišujeme několik hlavních paradigmat strojového učení, z~nichž každé má své typické úlohy a~metody:

\begin{enumerate}
    \item \textbf{Učení s~učitelem (Supervised Learning)} – modely se učí na označených datech, kde
    jsou vstupy i~správné výstupy známé, a~cílem je minimalizovat rozdíl mezi predikcemi modelu a~
skutečnými výstupy.
    \item \textbf{Učení bez učitele (Unsupervised Learning)} – model hledá strukturu nebo vzory v~
neoznačených datech bez explicitně definovaných výstupů.
    \item \textbf{Částečně řízené učení (Semi-Supervised Learning)} – kombinuje omezené množství
    označených dat s~velkým množstvím neoznačených dat pro efektivnější učení.
    \item \textbf{Zpětnovazební učení (Reinforcement Learning, RL)} – agent se učí sekvenci rozhodnutí
    pomocí interakcí s~prostředím, kdy dostává odměny podle kvality svých akcí \cite{Sutton2018}.
\end{enumerate}

V~této práci se zaměřujeme výhradně na zpětnovazební učení. RL se liší od ostatních paradigmat
především tím, že agent nepracuje pouze s~pevnými daty, ale aktivně ovlivňuje svět kolem sebe,
sleduje výsledky svých rozhodnutí a~učí se na základě toho. Tento způsob učení se formálně modeluje
pomocí \textit{Markovských rozhodovacích procesů (MDP)} – matematického rámce pro rozhodování v~
prostředí s~náhodnými přechody.

MDP je definován jako pětice $(\mathcal{S},\mathcal{A},P,R,\gamma)$, kde:

\begin{itemize}
    \item $\mathcal{S}$ je množina možných stavů prostředí,
    \item $\mathcal{A}$ je množina možných akcí, které může agent vykonat,
    \item $P(s' \mid s,a)$ je přechodová pravděpodobnost, tj. pravděpodobnost, že agent skončí ve stavu
    $s'$ poté, co ve stavu $s$ zvolí akci $a$,
    \item $R(s,a)$ je okamžitá odměna přidělená agentovi za provedení akce $a$ ve stavu $s$,
    \item $\gamma \in [0,1]$ je diskontní faktor, který určuje, jak moc si agent cení budoucích odměn
    vůči těm bezprostředním \cite{Sutton2018}.
\end{itemize}

Agent volí akce podle tzv. \textit{politiky} $\pi(a\mid s)$ – funkce, která pro každý stav $s$ přiřazuje
pravděpodobnost výběru akce $a$. RL algoritmy se snaží nalézt optimální politiku $\pi^*$, která maximalizuje
očekávanou sumu diskontovaných odměn za dlouhé časové horizonty.

\subsection{Proximal Policy Optimization (PPO)}

Proximal Policy Optimization (PPO) je jeden z~nejrozšířenějších algoritmů hlubokého zpětnovazebního učení,
který byl navržen výzkumníky OpenAI jako efektivní a~stabilní alternativa k~dřívějším policy gradient a~trust region
metodám \cite{PPOOpenAI}\cite{Schulman2017}\cite{Jochec2024}. PPO spojuje relativní jednoduchost implementace s~vysokou
stabilitou učení a~efektivním využitím nasbíraných dat.

Na rozdíl od hodnotových metod, jako je Q-learning \cite{Watkins1992}\cite{Sutton2018}, které se snaží
nejprve odhadnout hodnoty stavů či stav-akce a~teprve poté z~nich generovat politiku, PPO přímo optimalizuje
politiku parametrizovanou parametry $\theta$ (např. vahami neuronové sítě). Cílem je nalézt taková
parametry $\theta$, která vedou k~maximalizaci dlouhodobé očekávané odměny.

Klíčovým problémem u~jednodušších policy gradient metod je to, že jednotlivé aktualizace mohou
způsobit příliš velké změny v~politice, což může destabilizovat proces učení. Metody jako TRPO
(Trust Region Policy Optimization) řeší tento problém zavedením explicitních omezení na velikost
kroku při aktualizaci politiky, ale tyto přístupy jsou výpočetně náročné a~komplikované na implementaci
\cite{Schulman2015}. PPO tento problém řeší pomocí tzv. \textit{clipped surrogate objective}, která
udržuje aktualizace politiky v~kontrolovaném intervalu bez složitých optimalizačních podproblémů.

\newpage

Nejprve se zavádí poměr pravděpodobností mezi novou a~starou politikou:

\begin{equation}
r_t(\theta) = \frac{\pi_\theta(a_t\mid s_t)}{\pi_{\theta_{\text{old}}}(a_t\mid s_t)},
\end{equation}

kde $\pi_\theta$ je nová politika, $\pi_{\theta_{\text{old}}}$ je politika před aktualizací, $s_t$
je aktuální stav a $a_t$ akce zvolená ve stavu $s_t$. Tento poměr udává, jak moc se změnila
pravděpodobnost výběru konkrétní akce mezi starou a~novou politikou.

Poté se definuje \textit{oříznutý surrogate loss}:

\begin{equation}\label{eq:ppo_clip_latex}
L^{\text{CLIP}}(\theta) = \hat{\mathbb{E}}_t\left[\min\left(
r_t(\theta)\,\hat{A}_t,\ 
\operatorname{clip}\big(r_t(\theta),1-\epsilon,1+\epsilon\big)\,\hat{A}_t
\right)\right],
\end{equation}

kde:
\begin{itemize}
    \item $\hat{A}_t$ je odhad tzv. \textit{advantage funkce}, která udává, jak moc je akce $a_t$
    ve stavu $s_t$ lepší oproti průměrnému chování,
    \item $\epsilon$ je hyperparametr určující rozsah, ve kterém změna politiky není penalizována,
    \item $\operatorname{clip}$ ořízne hodnotu $r_t(\theta)$ do intervalu $[1-\epsilon, 1+\epsilon]$,
    čímž se minimalizují příliš velké změny politiky \cite{Schulman2017}.
\end{itemize}

Výraz \eqref{eq:ppo_clip_latex} říká, že pokud je změna politiky menší než povolené pásmo dané
$\epsilon$, ponechá se původní hodnota. Pokud je změna větší, bere se hodnota oříznutá na hranici
přijatelného rozsahu. Tento postup zabraňuje vykročení politiky do oblastí, kde by mohla být
nová politika nestabilní nebo výrazně horší než ta původní.

\newpage

Obrázek~\ref{obrPPO} vizualizuje architekturu PPO a~jeho interakci s~prostředím:

\begin{figure}[H]
\includegraphics[scale=4]{1-4.png} 
\centering
\caption{Vizualizace architektury algoritmu PPO. Agent je ovlivňován prostředím, zatímco dvě neuronové sítě — actor a~critic — slouží k~výběru akcí a~odhadu hodnoty stavu. Politika se aktualizuje pomocí oříznutého objektivu (clipped objective) a~výhody se počítají pomocí metody Generalized Advantage Estimation (GAE).\cite{PPO2021}}
\label{obrPPO} 
\end{figure}

PPO lze dále rozšířit o~dodatečné složky, které stabilizují učení a~podporují exploraci prostředí:

\begin{equation}
L^{\text{PPO}}(\theta) =
\hat{\mathbb{E}}_t\left[L^{\text{CLIP}}_t(\theta)-c_1\,L^{\text{VF}}_t(\theta)+c_2\,\mathcal{S}(\pi_\theta,s_t)\right],
\end{equation}

kde $L^{\text{VF}}$ je ztráta hodnotové funkce, $\mathcal{S}(\pi_\theta,s_t)$ entropie politiky ve
stavu $s_t$ a $c_1,c_2$ jsou váhové koeficienty určující jejich vliv na celkový loss \cite{Schulman2017}.

\newpage

\subsection{Deep Q-Network}

Algoritmus Deep Q-Network (DQN)\cite{Sutton2018}\cite{Mnih2013}\cite{Jochec2024}\cite{Mnih2015} představuje metodu zpětnovazebního učení, která kombinuje klasický Q-learning\cite{QLearning} s~hlubokými neuronovými sítěmi pro aproximaci akčně-hodnotové funkce \(Q(s,a;\theta)\). Tato funkce udává očekávanou kumulovanou diskontovanou odměnu, kterou agent získá, pokud ve stavu \(s\) provede akci \(a\) a~poté pokračuje optimálním chováním.

Teoretickým základem DQN je Bellmanova optimální rovnice:

\begin{equation}
Q^*(s_t,a_t)=\mathbb{E}\Big[r_{t+1}+\gamma\,\max_{a'}Q^*(s_{t+1},a')\;\Big|\;s_t,a_t\Big],
\end{equation}

kde \(s_t\) je aktuální stav, \(a_t\) zvolená akce, \(r_{t+1}\) okamžitá odměna, \(s_{t+1}\) následující stav a \(\gamma\in[0,1]\) diskontní faktor, který určuje relativní váhu budoucích odměn oproti bezprostředním.

V~klasickém Q-learningu se toto řeší pomocí aktualizačního pravidla:

\begin{equation}
Q(s_t,a_t)\leftarrow Q(s_t,a_t)+\alpha\Big(r_{t+1}+\gamma\,\max_{a'}Q(s_{t+1},a')-Q(s_t,a_t)\Big),
\end{equation}

kde \(\alpha\) je rychlost učení. Pro prostory se stovkami či tisíci stavů by taková tabulka byla nepraktická nebo nemožná. DQN proto využívá neuronovou síť parametrizovanou \(\theta\), která na vstup přijímá reprezentaci stavu \(s\) a~na výstup generuje \textbf{vektor Q-hodnot pro všechny akce}:

\begin{equation}
Q(s,a;\theta)\approx Q^*(s,a).
\end{equation}

Trénink DQN probíhá minimalizací \textbf{časového rozdílu (TD error)} mezi odhadovanou Q-hodnotou a~cílovou hodnotou odvozenou z~Bellmanovy rovnice. Kódování těchto odhadů a~stabilizace učení používá dvě hlavní techniky:

\begin{itemize}
    \item \textit{Přehrávací paměť} (\(\mathcal{D}\)) – ukládá přechody \((s_t,a_t,r_{t+1},s_{t+1})\) a~umožňuje náhodné vzorkování mini-batchů. Tím se přerušuje autokorelace mezi po sobě následujícími vzorky a~zvyšuje se rozmanitost tréninkových dat.
    \item \textit{Cílová síť} – označovaná \(Q(s,a;\theta^-)\), což je periodicky aktualizovaná kopie online sítě. Cílová síť poskytuje stabilní referenční hodnoty během učení.
\end{itemize}

\newpage

Pro výpočet cílové hodnoty a~ztrátové funkce DQN se používají následující výrazy.

Nejprve definujeme cílovou hodnotu \(y_t\):

\begin{equation}
y_t = r_{t+1} + \gamma\,\max_{a'}Q(s_{t+1},a';\theta^-).
\end{equation}

Tento výraz vyjadřuje tzv. \textit{temporal difference target}. Skládá se ze dvou částí:
\begin{itemize}
    \item \textit{Okamžitá odměna} \(r_{t+1}\): okamžitý zisk, který agent získá provedením akce \(a_t\) ve stavu \(s_t\).
    \item \textit{Odhad budoucích odměn} \(\gamma\,\max_{a'}Q(s_{t+1},a';\theta^-)\): diskontovaná hodnota nejlepší možné akce v~následujícím stavu \(s_{t+1}\), podle cílové sítě s~parametry \(\theta^-\). Tento člen reprezentuje odhad kumulativní budoucí odměny.
\end{itemize}

Hyperparametr \(\gamma\) je diskontní faktor, který určuje relativní váhu budoucích odměn vůči bezprostřední odměně. V~praxi se hodnoty \(\gamma\) blíží 1~v~úlohách, kde je důležitý dlouhodobý výsledek, a~menší \(\gamma\) se používá tam, kde jsou odměny krátkodobé.

Následuje výpočet ztrátové funkce \(L_t(\theta)\):

\begin{equation}
L_t(\theta)=\mathbb{E}_{(s,a,r,s')\sim\mathcal{D}}\Big[\big(y_t-Q(s_t,a_t;\theta)\big)^2\Big].
\end{equation}

Tato funkce vyjadřuje \textbf{střední kvadratickou chybu (mean squared error)} mezi:
\begin{itemize}
    \item \textit{aktuálním odhadem} hodnoty Q získaným z~online sítě \(Q(s_t,a_t;\theta)\), kterou trénujeme,
    \item \textit{cílovou hodnotou} \(y_t\), která využívá cílovou síť a~okamžitou odměnu.
\end{itemize}

Minimalizací této chyby učíme síť tak, aby její predikce Q-hodnot co nejlépe odpovídaly cílovým Bellmanovým odhadům. Důležité je, že cílové hodnoty se během tréninku *nepřenastavují okamžitě*, ale fixují se na delší interval pomocí zmíněné cílové sítě \(Q(s,a;\theta^-)\). To zásadně pomáhá stabilizovat učení a~zabraňuje oscilacím nebo divergenci při aktualizaci parametrů.

Po definici loss funkce následuje \textbf{aktualizace parametrů} pomocí gradientního sestupu:

\begin{equation}
\theta\leftarrow\theta-\alpha\,\nabla_{\theta}L_t(\theta),
\end{equation}

kde:
\begin{itemize}
    \item \(\alpha\) je rychlost učení (learning rate), která určuje, jak velký krok do směru negativního gradientu uděláme v~každé aktualizaci,
    \item \(\nabla_{\theta}L_t(\theta)\) je gradient loss funkce vůči parametrům sítě \(\theta\). Tento gradient udává, jak změna každého parametru ovlivní chybu.
\end{itemize}

Gradientní sestup směřuje parametry sítě tak, aby se predikovaná Q-hodnota \(Q(s_t,a_t;\theta)\) blížila cílovému odhadu \(y_t\). V~důsledku toho síť postupně zlepšuje své odhady Q-hodnot pro různé stav-akce páry.

Poslední rovnice vyjadřuje \textbf{synchronizaci parametrů cílové sítě} s~hlavní sítí:

\begin{equation}
\theta^-\leftarrow\theta.
\end{equation}

Tento krok se neprovádí při každé aktualizaci, ale periodicky (např. každých \(K\) tréninkových kroků). V~praxi se tím zajistí, že cílová síť zůstane dostatečně „fixní“ po určitou dobu, což zabraňuje tomu, aby se cílové hodnoty měnily příliš rychle a~destabilizovaly tréninkový proces. Bez této separace by totiž síť během učení upravovala své parametry tak, že by se sama učila na neustále se měnící cílové funkci — což má tendenci vést k~oscilacím nebo divergenci.

Stručně řečeno:
\begin{itemize}
    \item \textit{Loss funkce} měří, jak moc se predikce odchyluje od cílové hodnoty na základě Bellmanovy rovnice.
    \item \textit{Gradientní aktualizace} upravuje váhy tak, aby se tato odchylka minimalizovala.
    \item \textit{Target network} odděluje výpočet cílových Q-hodnot od parametrů, které se právě učí, čímž stabilizuje trénink.
\end{itemize}

Tréninkový cyklus DQN tedy typicky probíhá následovně:

\begin{enumerate}
    \item Sledování aktuálního stavu \(s_t\).
    \item Výběr akce \(a_t\) pomocí \(\epsilon\)-greedy politiky.
    \item Provedení akce, získání odměny \(r_{t+1}\) a~nového stavu \(s_{t+1}\).
    \item Uložení přechodu do paměti \(\mathcal{D}\).
    \item Náhodné vzorkování mini-batchů dat z \(\mathcal{D}\).
    \item Výpočet cílové hodnoty \(y_t\) a~minimalizace TD chyby.
    \item Periodická aktualizace parametrů cílové sítě (\(\theta^- \leftarrow \theta\)).
\end{enumerate}

Algoritmus DQN je citlivý na nestabilitu a~divergenci kvůli kombinaci nelineární aproximace neuronovou sítí a~bootstrapovaných aktualizací. Přehrávací paměť redukuje autokorelaci dat a~zvyšuje efektivitu využití zkušeností, zatímco cílová síť poskytuje stabilnější referenční hodnoty. Pokročilé varianty, jako jsou \textit{Double DQN}, \textit{Dueling DQN} nebo \textit{Prioritized Experience Replay}, rozšiřují základní metodu s~cílem zlepšit konvergenci, stabilitu a~efektivitu učení v~komplexních prostředích.

\begin{figure}[ht]
\includegraphics[scale=0.45]{1-7.png} 
\centering
\caption{Příklad architektury Deep Q-Learning\cite{Elsayed2021}}
\label{obrADQL}
\end{figure}

Tyto dva algoritmy byly již zmíněny ve výzkumné práci\cite{Jochec2024}, kde tato diplomová práce na ni navazuje. Teď budou zmíněny dvě trénovací techniky, Self-Play a~Population Based Training.

\newpage

\subsection{Self-Play}

Self-play\cite{Zhang2024} je technika tréninku v~oblasti zpětnovazebního učení, která umožňuje agentům 
zlepšovat své strategie prostřednictvím interakcí s~jinými instancemi sebe sama – typicky s~kopiemi, 
historickými verzemi nebo současnými protivníky založenými na stejné struktuře politiky. Cílem self-play 
je generovat adaptivní tréninková data bez nutnosti externích protivníků či lidského dozoru, čímž se agent 
učí robustní strategii na základě postupného zvyšování úrovně soupeře během tréninku.

Na rozdíl od klasického zpětnovazebního učení, kde agent interaguje se statickým prostředím, self-play funguje 
v~rámci dynamického prostředí, které vytváří další agentní entity, jejichž chování se v~čase vyvíjí a~
reaguje na aktuální strategii. To vede k~učením, která mohou překročit jednoduché optimalizační úlohy a~
dosahovat hlubší strategické adaptace, zejména v~úlohách s~konkurenční povahou.

Základní forma self-play, často označovaná jako \uv{vanilla self-play}, spočívá v~tom, že agent opakovaně 
hraje proti své vlastní nejnovější verzi. Tato metoda funguje efektivně v~prostředích, kde lze strategie 
řadit podle jejich kvality (transitivní hry), protože agent postupně zvyšuje úroveň soupeře sám proti sobě. 
V~prostředích s~ne-transitivní dynamikou může tato jednoduchá forma self-play vést k~cyklickému chování 
a~suboptimálním strategiím.

\newpage

Aby se tyto nevýhody zmírnily, vznikají různé varianty:

\begin{itemize}
    \item \textit{Fiktivní self-play} rozšiřuje trénink tím, že agent hraje proti průměru všech svých 
    předchozích verzí politiky. Tím se zvyšuje robustnost učení a~snižuje riziko zacyklení v~suboptimálních strategiích.
    \item \textit{Prioritizovaná fiktivní self-play} používá vážený výběr soupeřů, kde se preferují ty 
    politiky, proti kterým si agent vede méně dobře – tedy trénuje intenzivněji vůči slabým stránkám vlastní strategie.
    \item \textit{$\delta$-uniform self-play} zavádí hyperparametr $\delta$, který určuje, jaká část 
    historické populace politik se použije jako soupeři. Nižší $\delta$ zahrnuje širší spektrum historických politik, 
    zatímco vyšší $\delta$ dává přednost nedávno generovaným verzím.
\end{itemize}

Self-play se často integruje s~tzv. \textit{population-based} přístupy. Namísto jednoho agenta existuje populace 
politik $\Pi$, která se v~čase rozšiřuje o~nové verze. Každá politika se trénuje vůči soupeřům vybraným z~populace 
na základě určité metastrategie $\Sigma$, často reprezentované interakční maticí nebo distribuční strategií. Tento 
proces vytváří bohatý a~adaptivní prostor pro evoluci strategií a~vzájemnou adaptaci politik.

Základní smyčka self-play s~populací obvykle probíhá následovně:

\begin{enumerate}
    \item Nové politiky jsou inicializovány, často na základě posledních iterací existujících politik.
    \item Každá politika se trénuje proti soupeřům vybraným podle metastrategie $\Sigma$.
    \item Nově natrénované politiky se přidají zpět do populace a~proces se opakuje.
\end{enumerate}

Z~hlediska výpočetní efektivity se rozlišují varianty jako:

\begin{itemize}
    \item \textit{lazy initialization}, kdy se politiky inicializují až v~okamžiku potřeby,
    \item \textit{immediate initialization}, kdy se populace vytvoří celá najednou na začátku.
\end{itemize}

Dále se rozlišuje:

\begin{itemize}
    \item \textit{standardní trénink}, kde se v~každé iteraci trénuje jeden agent,
    \item \textit{průběžný trénink}, kdy se všechny politiky trénují paralelně.
\end{itemize}

Zejména v~kontextu větších populací a~her s~komplexní strategií hraje důležitou roli tzv. PSRO (\textit{Policy-Space Response Oracles}). 
PSRO rozšiřuje klasické self-play tak, že pracuje s~populací politik a~iterativně generuje nejlepší odpovědi na směs politik populace, 
čímž se přibližuje rovnováze typu Nashovy rovnováhy. To je důležité zejména v~nestranově tranzitivních nebo multi-agentních hrách, 
kde jednoduché sebehráčské strategie nemusí konvergovat ke stabilnímu optimu.

\newpage

Self-play nachází využití v~konkurenčních a~multi-agentních prostředích — klasickými příklady jsou hry jako Go, šachy, StarCraft II 
nebo Dota 2 — a~také v~jiných oblastech, například v~některých aplikacích robotiky, kde dynamické protinavádění podporuje robustní 
strategie bez nutnosti lidského tréninkového dozoru.

Obrázek~\ref{obrSPvsPPvsBCP} ilustruje rozdíly mezi přístupy \uv{self-play} (SP), \uv{population-play} (PP) a \uv{behavioral cloning play} (BCP). 
Druhý ze zmíněných přístupů využívá širší populace politik, zatímco BCP je založen na učení z~lidských demonstračních dat, což je odlišné od 
čistě self-play přístupů, které generují data prostřednictvím interakcí agentů mezi sebou.

\begin{figure}[ht]
\includegraphics[scale=0.67]{1-5.png} 
\centering
\caption{Srovnání tří přístupů k~tréninku agentů: \uv{self-play} (SP), kde agent hraje sám proti sobě; \uv{population-play} (PP), kde interaguje s~různými agenty v~populaci; a \uv{behavioral cloning play} (BCP), kde se učení zakládá na lidských demonstračních datech. Upraveno podle zdroje\cite{Strouse2021}}
\label{obrSPvsPPvsBCP} 
\end{figure}

\newpage

\subsection{Population Based Training}

Population-Based Training (PBT)\cite{Jaderberg2017} je technika automatického ladění hyperparametrů, která propojuje samotný trénink neuronových sítí s~jejich optimalizací tak, aby se váhy sítě a~její hyperparametry upravovaly současně během jednoho běhu učení. PBT vzniká jako reakce na omezení tradičních metod ladění, jako je grid search nebo Bayesian optimization, které často vyžadují mnoho samostatných tréninků s~fixními konfiguracemi a~mohou být časově i~výpočetně náročné.

V~PBT se nevyhledává pouze jedna nejlepší sada hyperparametrů, ale algoritmus se snaží \textbf{automaticky objevit jejich optimální plán (schedule)} během tréninku samotného modelu — díky tomu se může přizpůsobovat různým fázím učení, kdy různé hodnoty hyperparametrů jsou vhodné v~různých časech tréninku.

Na začátku PBT inicializuje populaci modelů, přičemž každá instance populace má:
\begin{itemize}
    \item vlastní váhy (parametry neuronové sítě),
    \item vlastní konfiguraci hyperparametrů (např. learning rate, entropická penalizace, délka unrollu apod.).
\end{itemize}

Každý model se trénuje paralelně a~nezávisle pomocí běžných optimalizačních metod (např. SGD nebo Adam). V~pravidelných intervalech se populace vyhodnocuje a~provádí tzv. \textit{exploit-and-explore} fáze, které spojují principy evolučních algoritmů s~gradientním učením.

\textit{Exploit} znamená, že modely s~horším výkonem přebírají váhy a/nebo hyperparametry od modelů s~lepším výkonem. Tato operace je analogická selekci v~genetických algoritmech — silnější jedinci „přežijí“ a~jejich parametry se šíří do dalších populací. 
\textit{Explore} znamená, že tyto převzaté hyperparametry jsou \textbf{lehce modifikovány} (např. vynásobením faktorem 1{,}2~nebo 0{,}8), aby populace prozkoumala nové konfigurace. Tím se neustále testují různé kombinace nastavení během tréninku a~umožňuje se adaptivní ladění hyperparametrů bez nutnosti opakovaných běhů tréninku od začátku.

Proces lze shrnout následovně:
\begin{enumerate}
    \item Každý model ve své instanci trénuje standardní metodou s~aktuální konfigurací hyperparametrů.
    \item Po určité době nebo podle metriky se vyhodnocuje výkon každého modelu.
    \item Pokud je model označen jako \uv{ready}, probíhá fáze \uv{exploit} – model si kopíruje váhy a~hyperparametry od výkonnějších modelů v~populaci.
    \item Následuje fáze \uv{explore}, při které se hyperparametry lehce mění (například násobky, malé perturbace nebo náhodné přegenerování).
    \item Trénink pokračuje s~nově nastavenými vahami a~hyperparametry a~celý cyklus se opakuje.
\end{enumerate}

Tento způsob optimalizace kombinuje směry gradientního učení a~evolučního vyhledávání tak, že:
\begin{itemize}
    \item zajišťuje \textbf{adaptivní změnu hyperparametrů během učení},
    \item eliminuje potřebu řešit tréninky s~fixními konfiguracemi opakovaně,
    \item dovoluje paralelní škálování na heterogenních výpočetních systémech bez centrální koordinace.
\end{itemize}

Obrázek~\ref{obrPBT} ilustruje, jak si modely v~rámci populace sdílejí váhy a~mění hyperparametry na základě výkonnosti:

\begin{figure}[ht]
\includegraphics[scale=0.35]{1-6.png} 
\centering
\caption{Population-Based Training neuronových sítí začíná jako náhodné prohledávání, ale v~průběhu tréninku umožňuje jednotlivým modelům využít dílčí výsledky ostatních (fáze \uv{exploit}) a~zároveň zkoušet nové kombinace hyperparametrů (fáze \uv{explore}).\cite{PBT2017} }
\label{obrPBT} 
\end{figure}

Klíčové principy PBT mohou být interpretovány jako hybridní systém:
\begin{itemize}
    \item \textit{výběr modelů} podobný truncation nebo turnajové selekci v~evolučních metodách,
    \item \textit{perturbace hyperparametrů} jako jednoduchý heuristický způsob prozkoumání prostoru hyperparametrů během běhu tréninku,
    \item \textit{asynchronní aktualizace}, která eliminuje nutnost synchronizace všech modelů.
\end{itemize}

Tato technika byla v~původní práci demonstrována na úlohách hlubokého zpětnovazebního učení, kde vedla nejen k~rychlejší konvergenci, ale i~k~vyššímu konečnému výkonu modelů při optimalizaci různých hyperparametrů. PBT také nachází uplatnění v~dalších oblastech, jako je strojové překládání nebo trénink generativních modelů, kde automaticky objevuje sekvence hyperparametrů, které překonávají ručně laděné konfigurace.

Důležitou výhodou PBT je jeho \textbf{online charakter} — ladění hyperparametrů probíhá současně s~tréninkem modelů, což je efektivnější nejen časově, ale i~z~hlediska výpočetních zdrojů, protože se eliminuje potřeba trénovat modely od začátku po každé úpravě hyperparametrů. Tento aspekt výrazně odliší PBT od tradičních metod ladění, které obvykle vyžadují dokončení jednoho tréninkového běhu před dalším experimentem.

Celkově lze PBT vnímat jako techniku, která do procesu učení vnáší adaptivitu, sdílení informací mezi modely a~evoluční principy, což podporuje robustnější a~efektivnější optimalizaci ve srovnání s~pevně danými metodami ladění.

\newpage

\section{AI agenti ve hrách}
\subsection{AlphaStar}

AlphaStar\cite{Vinyalsbr} je agent umělé inteligence (AI) vyvíjený společností DeepMind pro hraní hry StarCraft II\cite{WikiSC}, což je komplexní a~dynamická online strategie v~reálném čase s~částečnou pozorovatelností, rozsáhlým akčním prostorem a~dlouhodobým plánováním. Hra StarCraft II je považována za jeden z~nejnáročnějších testovacích prostorů pro vývoj inteligentních systémů, protože kombinuje strategické a~taktické rozhodování se zpracováním neurčitých informací a~interakcí více agentů.

AlphaStar byla veřejnosti představena v~lednu 2019 a~již v~srpnu téhož roku dosáhla hodnosti \uv{Grandmaster}, což je nejvyšší úroveň na oficiální žebříčkové lize StarCraft II a~znamená, že se agent řadí mezi přibližně 0{,}2\,\% nejlepších lidských hráčů. Tento úspěch představuje významný milník nejen v~oblasti herní AI, ale i~ve vývoji komplexních autonomních systémů, které zvládají vysokodimenzionální rozhodovací úlohy.

StarCraft II jako prostředí klade na agenta řadu výzev, které nejsou přítomné v~jednodušších hrách nebo simulacích: \textbf{časové omezení, skryté informace (fog of war), rozsáhlý akční prostor, dynamické kombinace jednotek a~strategií, a~nutnost plánování na krátkodobé i~dlouhodobé úrovni}. Teprve kombinace moderních technik strojového učení dokázala tyto výzvy překonat.

Vývoj AlphaStar odráží dlouhodobý směr výzkumu DeepMind. Zakladatel Demis Hassabis již v~roce 2011 označil StarCraft za další krok ve výzkumu herní AI po úspěších s~AlphaGo, agentem pro hru Go, který porazil světového šampiona. Po spojení s~Blizzard Entertainment a~zpřístupnění API pro boty v~roce 2017 začal DeepMind zpracovávat historická herní data a~postupně stavět a~trénovat své agentní sítě.

AlphaStar funguje na kombinaci několika technik:
\begin{itemize}
    \item \textit{Učení z~lidských dat}: původní agent byl trénován pomocí imitace strategií z~rozsáhlé databáze anonymních her lidských hráčů, aby se naučil základní struktury herního rozhodování.
    \item \textit{Hluboké zpětnovazební učení (Reinforcement Learning)}: agent dále optimalizoval svou politiku pomocí zpětné vazby ze simulovaných her tak, aby maximalizoval pravděpodobnost vítězství.
    \item \textit{Multi-agentní trénink a~self-play}: agenti hráli proti sobě v~rozsáhlých tréninkových ligách, což umožňovalo objevování a~adaptaci různých strategií bez ručního navrhování pravidel.
    \item \textit{Population-Based Training (PBT)}\cite{Jaderberg2017}: tato technika automaticky ladila hyperparametry během tréninku, čímž zvyšovala efektivitu a~kvalitu učení.
\end{itemize}

AlphaStar je navržen tak, aby zvládal komplexní rozhodovací prostor StarCraft II pomocí hlubokých neuronových sítí, které zpracovávají game-state v~reálném čase a~generují akce, které reflektují kombinaci taktického i~strategického myšlení.

AlphaStar se stal historicky prvním AI systemem, který dosáhl \textbf{hodnosti Grandmaster ve veřejném žebříčku hry StarCraft II} bez umělých omezení algoritmu, čímž překonal většinu lidských hráčů v~rámci oficiálního kompetitivního prostředí. Tento výkon byl dosažen i~přesto, že prostředí kladlo na agenta lidské omezení jako omezené zorné pole kamery a~omezenou rychlost akcí, aby se co nejvíce přiblížilo podmínkám, za kterých hrají lidé.

Model AlphaStar tak představuje \textbf{pokročilou kombinaci imitace, reinforcement learningu, multi-agentního tréninku a~adaptivního ladění}, která umožňuje extrémně komplexní herní rozhodování. Přestože je tento systém navržen primárně pro herní prostředí, principy, které byly při jeho vývoji použity, mají potenciál být aplikovány i~v~jiných oblastech, kde je potřeba řešit nelineární rozhodovací úlohy v~přítomnosti neúplné informace a~dynamických protivníků.

AlphaStar nepředstavuje jedinou strategii, ale udržuje populaci různorodých politik, které odrážejí různé herní styly a~strategie. Tato populace se během tréninku dynamicky rozšiřuje a~přizpůsobuje, čímž se zvyšuje robustnost a~schopnost adaptace agentů v~různorodém prostředí. Tento přístup je úzce spojen s~principy kvalitní diverzity (Quality Diversity, QD), kdy místo jediné optimalizované strategie vzniká více kvalitních, ale odlišných řešení, která jsou vůči sobě navzájem méně zranitelná.

\begin{figure}[H]
\includegraphics[scale=0.41]{1-1.jpg}
\centering
\caption{Vizualizace rozhodovacího procesu AlphaStaru během zápasu proti hráči. Systém zpracovává syrová pozorování, aktivuje neuronové sítě, vyhodnocuje potenciální akce a~předpovídá výsledky zápasu.\cite{AlphaStar2019}}
\label{obrAlphaStar}
\end{figure}

\newpage

\subsection{OpenAI Five}

OpenAI Five\cite{OpenAIFive} je agent umělé inteligence vyvinutý společností OpenAI, navržený pro hraní týmové videohry Dota 2\cite{Dota2}, ve které proti sobě nastupují dvě pětičlenná družstva. Tato hra představuje jedno z~nejkomplexnějších prostředí pro testování agentů umělé inteligence, protože kombinuje \textbf{dlouhé časové horizonty, částečnou pozorovatelnost a~vysokodimenzionální akční a~stavový prosto¨r}, které výrazně přesahují jednodušší hry jako šachy či Go.

Poprvé se OpenAI Five veřejnosti představuje v~roce 2017 během živé demonstrace na turnaji The International, kde v~duelu jeden na jednoho poráží profesionálního hráče Dendiho — což ukazuje, že agent dokáže získat přehled a~efektivně rozhodovat i~v~takto náročném prostředí. Následující rok systém pokročuje natolik, že zvládá ovládat celý tým pěti agentů a~poráží amatérské i~poloprofesionální lidské soupeře ve formátu 5v5.

Vývoj agentních algoritmů pro OpenAI Five začíná již v~listopadu 2016. OpenAI zvolila Dota 2~jako testovací prostředí mimo jiné kvůli její \textbf{složitosti, nepředvídatelnosti hry, robustnímu API a~vysoké popularitě}, což umožnilo testovat a~měřit výkon agentů v~realistickém konkurenčním scénáři.

OpenAI Five využívá \textbf{framework hlubokého zpětnovazebního učení (deep reinforcement learning)}, přičemž jádrem tréninkové metodiky je algoritmus \uv{Proximal Policy Optimization} (PPO)\cite{PPOOpenAI}\cite{Schulman2017}. PPO je gradientní algoritmus politiky, který kombinuje efektivní využití dat, stabilní učení a~relativně jednoduchou implementaci — což je důležité při škálování učení na rozsáhlé distribuované systémy, jaké OpenAI v~tomto projektu použila. PPO zabraňuje příliš razantním aktualizacím politiky zavedením tzv. oříznutého objektivu (\textit{clipped surrogate objective}), který omezuje změny v~distribučních pravděpodobnostech akcí během gradientních kroků.

Každý agent v~OpenAI Five je reprezentován neuronovou sítí s~jednou vrstvou typu \textbf{LSTM o~4096 jednotkách}, která zpracovává aktuální herní stav získaný přes API Dota 2. Tento stav je zakódován jako vektor zhruba 20 000 číselných hodnot, jež popisují relevantní informace z~herního prostředí, například pozice jednotek, dostupné akce, cooldowny schopností a~další proměnné.

Rozhodování probíhá pomocí několika \textbf{nezávislých akčních hlav}, z~nichž každá odpovídá za jiný aspekt chování agenta — například volbu konkrétní akce, její umístění v~prostoru nebo časování provedení akce. Každý krok agenta je reprezentován osmi číselnými hodnotami (tzv. \uv{action enumeration}), které definují konkrétní volbu v~daném časovém okamžiku.

Trénink OpenAI Five probíhá na rozsáhlé distribuované infrastruktuře nazvané Rapid, která umožňuje masivní paralelní výpočet a~koordinaci tisíců strojů. Díky tomu systém odehrává denně obrovské množství simulovaných her — v~některých fázích tréninku se sbírá ekvivalent více než 180 let herního času na den, což je umožněno paralelním hraním \uv{self-play} na stovkách GPU a~desítkách tisíc CPU jader.

\newpage

OpenAI Five tak neustále trénuje a~optimalizuje své politiky hraním proti sobě samému (\uv{self-play}), přičemž model získává odměny za různé klíčové herní cíle, jako je zabíjení nepřátel, ničení struktur nebo dosahování herních objektivů. Tento proces podporuje postupné zlepšování strategie a~koordinace celého týmu agentů.

Závěrečnou veřejnou prezentací projektu byly exhibiční zápasy v~dubnu 2019, během nichž OpenAI Five poráží tým OG — tehdejší mistry světa z~The International 2018 — ve formátu best-of-three se skóre 2:0~na živém streamu, čímž demonstruje, že jeho výkon překonává i~špičkové lidské profesionály. Během čtyřdenní online akce OpenAI Five dále otestoval své schopnosti v~více než 7~000 zápasech proti veřejným hráčům, z~nichž vyhrál přibližně 99{,}4 \% — což potvrzuje konzistentní a~robustní výkon v~rozsáhlém měřítku.

Díky své robustnosti, efektivitě a~schopnosti fungovat v~distribuovaných systémech je PPO ideálně vhodný pro tuto aplikaci. Umožňuje agentům provádět rozsáhlý paralelní trénink bez přeučení, přičemž jeho kompatibilita se sdílenými politikami i~hodnotovými funkcemi přispívá k~úspěšné koordinaci agentů v~komplexním prostředí Dota 2.

\newpage

\subsection{Voyager}

Voyager\cite{Wangbr} je zástupce tzv. ztělesněných agentů (embodied agents), kteří kombinují výhody velkých jazykových modelů (LLM) s~učením ve složitých, otevřených prostředích. Tento agent funguje autonomně ve hře Minecraft — bez přímého lidského zásahu nebo předem definovaných odměnových funkcí — a~učí se postupně z~vlastních interakcí s~herním světem. Minecraft je pro tento typ agenta ideálním prostředím, protože neobsahuje pevný konec hry, má procedurálně generované světy a~vyžaduje dlouhodobé plánování a~adaptivní rozhodování, což odráží výzvy reálných úloh.

Klíčové vlastnosti Voyagera jsou založeny na třech hlavních mechanismech, které mu umožňují \textbf{otevřené, kontinuální učení bez lidské intervence}:

\begin{enumerate}
    \item \textbf{Automatické kurikulum} – agent si sám generuje sekvenci úkolů na základě aktuálního stavu světa, inventáře a~historie již provedených akcí. Tato posloupnost úkolů maximalizuje exploraci a~postupně vede ke schopnostem potřebným pro složitější dovednosti.
    \item \textbf{Knihovna dovedností} – po úspěšném splnění úkolu agent vytvoří program reprezentující tuto dovednost a~uloží jej do knihovny. Dovednosti jsou indexovány vektorovými embeddingy jejich popisu, což umožňuje jejich pozdější opětovné použití nebo kombinaci s~jinými schopnostmi.
    \item \textbf{Iterativní promptování} – při řešení úkolu agent nejprve navrhne program (např. v~JavaScriptu) pro daný úkol pomocí GPT-4~a~poté jej provede ve hře. Výsledek, včetně chyb a~zpětné vazby z~prostředí, se vrací GPT-4, které program upraví a~iterativně opravuje, dokud není úkol úspěšně dokončen nebo vyhodnocen jako neřešitelný.
\end{enumerate}

Automatické kurikulum nehledá náhodné úkoly, ale generuje je kontextově na základě aktuální situace v~prostředí, což umožňuje efektivní prozkoumávání světa a~dlouhodobé zlepšování agentovy schopnosti učit se.

Dovednosti, které Voyager získá, nejsou izolované akce, ale mohou být hierarchicky složeny. Například dovednost vytvořit železný nástroj může využívat dovednosti jako těžba rudy, výroba pece nebo sběr surovin, přičemž se buduje na znalostech získaných dříve a~minimalizuje hrazené opakování již naučených postupů.

Iterativní promptování je charakteristické tím, že agent nevykonává pouze jednorázové akce, ale \textbf{opakuje cyklus generování a~hodnocení kódu} za použití zpětné vazby z~prostředí. Kód, který selže kvůli chybě, je upraven na základě analýzy chyby a~dalšího promptu, čímž vzniká robustní smyčka učení, která agentovi umožňuje zlepšovat své strategie bez explicitního překalibrování modelu.

Ve srovnávacích experimentech Voyager významně překonává předchozí agentní metody založené na jiných přístupech, jako jsou ReAct či Auto-GPT. Agent získává více unikátních předmětů, pokrývá větší území a~rychleji dosahuje technologických milníků, což ukazuje jeho schopnost adaptovat se a~učit se v~otevřeném prostředí bez ručního ladění cílů.

Celkově lze Voyagera vnímat jako významný krok směrem k~obecné autonomní umělé inteligenci v~otevřených světech, protože demonstruje, že \textbf{agent vybavený LLM a~mechanismy samostatného učení může dlouhodobě akumulovat znalosti, generalizovat získané dovednosti a~přizpůsobovat se novým úlohám bez lidského zásahu}.

\begin{figure}[ht]
\includegraphics[scale=1]{1-3.png} 
\centering
\caption{Vizualizace hlavních komponent systému Voyager – automatického kurikula, knihovny dovedností a~iterativního promptování.\cite{Voyager2023}}
\label{obrVoyager} 
\end{figure}